\section{Conclusions} \label{sec:concl}

The automatic recognition of the human group behavior is a challenging task inside the computer vision area due to the different aspects that define a behavior (e.g. movement, context, etc.). To address it, in this paper a multi-stream generic computer vision architecture model has been proposed. The proposed model is based on a movement descriptor (Deep-ADV) and a (deep learning) classification stage. By using a motion descriptor, the system generality increases regardless of the datasets as the classifier is trained with the descriptor features and not from the raw images. Moreover, smaller dataset is required to train the architecture because the solution space is reduced. The generic architecture model has been instantiated to classify multiple group behaviors (D-ADV-MC) and to detect anomalies (D-ADV-OC).

The D-ADV-MC instance uses a late fusion based two-stream architecture to classify different group activities. The performance of this instance has been evaluated using different datasets showing very high accuracy rate, outperforming current state-of-the-art proposals. Furthermore, the D-ADV-OC uses a three-stream architecture, one per movement component and a third one representing context. This instance has been evaluated using different public datasets with results in accordance with other state-of-the-art methods. In both instances, the use of transfer learning with models trained on Imagenet has been key to train the models with small datasets.

Future work will focus on the proposal of new configurations to calculate the D-ADV as well as the use of other deep architectures. In addition, the use of other context inputs (e.g., location, time, etc.) will be analyzed in detail.
